\documentclass[11pt,a4paper]{article}

\usepackage[T2A]{fontenc}
\usepackage[utf8]{inputenc}
\usepackage[russian,english]{babel}
\usepackage{amsmath,amssymb,amsthm}
\usepackage{geometry}
\usepackage{graphicx}
\usepackage{booktabs}
\usepackage{hyperref}
\usepackage{enumitem}
\usepackage{array}
\usepackage{multirow}
\usepackage{xcolor}
\usepackage{float}

\geometry{margin=1in}

\title{Angel-in-Pocket: субъектно-центрированная мультеверс-архитектура ИИ для личного сопровождения, обнуления шума и иерархической устойчивости}
\author{qqewq}
\date{\today}

\begin{document}

\maketitle

\begin{abstract}
В работе предлагается архитектура Angel-in-Pocket — субъектно-центрированная система искусственного интеллекта, предназначенная для персонального сопровождения одного конкретного пользователя. В отличие от универсальных ассистентов и многoагентных рыночных платформ, система организована вокруг одного субъекта как центральной точки отсчёта. Архитектура объединяет извлечение данных, ранжирование, байесовский учёт неопределённости, обнуление шума и паразитарных сигналов, иерархическую оценку устойчивости и долгосрочное прогнозирование. Мы вводим понятие личного \emph{мультеверта}: пространства решений, в котором все рекомендации интерпретируются относительно данных, целей, ограничений, здоровья и поведенческого контекста пользователя. Предлагаемая система ориентирована на практические применения в мониторинге здоровья, помощи при покупках, выборе услуг, оптимизации образа жизни и сценарном прогнозировании. В статье приведены формальные определения, архитектурные модули, цели оптимизации и пути реализации.
\end{abstract}

\section{Введение}

Большие языковые модели всё чаще используются как ассистенты, однако большинство существующих систем остаются реактивными, универсальными и слабо персонализированными. Angel-in-Pocket предлагает иной подход: постоянного, индивидуального и субъектно-центрированного ИИ-компаньона, который знает пользователя через явные данные, исторический контекст и модель предпочтений. Центральная идея состоит не в «мультиверсе» равноправных агентов, а в \emph{мультеверте}: персональном пространстве решений, организованном вокруг одного конкретного субъекта.

В данной статье формализуется эта концепция и её связь с семейством GRA-методов, включая иерархическую устойчивость, байесовское ранжирование, swarm-агрегацию, резонансное обнуление и переключение, ориентированное на субъекта. Предлагаемая архитектура не является медицинским или юридическим авторитетом; это система рекомендаций и прогнозирования, помогающая субъекту принимать более качественные решения.

\section{Мотивация}

Современный пользователь сталкивается с информационной перегрузкой в рынках, сервисах и системах здравоохранения. При покупках необходимо сравнивать цены, условия доставки, качество и скрытые риски. В медицине — интерпретировать анализы, следить за динамикой и решать, когда обращаться к специалисту. При планировании жизни — балансировать сон, работу, питание, стресс и долгосрочные цели.

Субъектно-центрированный ассистент может снизить эту нагрузку, выступая в роли постоянного фильтра решений. Он не просто отвечает на вопросы, а постоянно моделирует состояние пользователя, отбрасывает шум, оценивает неопределённость и предлагает действия, устойчивые как локально, так и в долгосрочной перспективе.

\section{Постановка задачи}

Пусть пользователь-субъект обозначен через $u$. В момент времени $t$ состояние субъекта задаётся вектором
\[
x_t = (h_t, b_t, e_t, s_t, p_t),
\]
где $h_t$ — состояние здоровья, $b_t$ — поведенческое состояние, $e_t$ — экономическое состояние, $s_t$ — состояние стресса, $p_t$ — состояние предпочтений.

Пусть имеется множество кандидатных действий или вариантов:
\[
A = \{a_1, a_2, \dots, a_n\}.
\]
Система ищет оптимальное действие $a^*$, максимизирующее полезность и минимизирующее риск и шум:
\[
a^* = \arg\max_{a_i \in A} \left[ U(a_i, u, x_t) - \lambda R(a_i, u, x_t) - \mu N(a_i) \right].
\]
Здесь $U$ — полезность, $R$ — риск, $N$ — штраф за обнуление, отражающий паразитарные или вводящие в заблуждение сигналы, а $\lambda, \mu \ge 0$ — параметры настройки.

\section{Архитектура}

Система организована в пять слоёв.

\subsection{Субъектное ядро}
Субъектное ядро хранит пользовательские атрибуты:
\begin{itemize}[leftmargin=1.2em]
    \item демографический профиль,
    \item привычки и распорядок,
    \item бюджетные ограничения,
    \item история здоровья,
    \item предпочтения и антипатии,
    \item цели и приоритеты.
\end{itemize}

\subsection{Слой извлечения данных}
Слой извлечения собирает кандидатные варианты из доступных источников:
\begin{itemize}[leftmargin=1.2em]
    \item торговые площадки,
    \item каталоги услуг,
    \item медицинские записи,
    \item загруженные документы,
    \item пользовательские заметки,
    \item внешние API или локальные базы данных.
\end{itemize}

\subsection{Слой обнуления}
Слой обнуления удаляет малодоверительные или шумные варианты. Он применяет резонансную фильтрацию, хиральную фильтрацию, оценку доверия к источнику и штрафы за несовпадение с субъектом. Формально:
\[
N(a_i) =
\begin{cases}
1, & \text{если } a_i \text{ проходит фильтры}, \\
0, & \text{иначе}.
\end{cases}
\]
Только кандидаты с $N(a_i)=1$ передаются на следующий этап.

\subsection{Слой ранжирования и устойчивости}
Каждый кандидат оценивается по взвешенной функции полезности:
\[
S_i = w_p p_i + w_d d_i - w_q q_i + w_r r_i - w_s s_i - w_h h_i + w_c c_i,
\]
где:
\begin{itemize}[leftmargin=1.2em]
    \item $p_i$ — цена,
    \item $d_i$ — стоимость или задержка доставки,
    \item $q_i$ — качество,
    \item $r_i$ — риск,
    \item $s_i$ — устойчивость,
    \item $h_i$ — прямая полезность,
    \item $c_i$ — совместимость с состоянием субъекта.
\end{itemize}

Выбор производится по минимальному эффективному скору:
\[
a^* = \arg\min_{a_i \in A,\, N(a_i)=1} \tilde{S}_i,
\]
где
\[
\tilde{S}_i = \mathbb{E}[S_i] - \alpha \mathrm{Stab}(a_i).
\]

\subsection{Слой оркестрации}
Слой оркестрации объединяет несколько моделей или эвристик через swarm-агрегацию:
\[
\hat{S}_i = \frac{1}{M} \sum_{m=1}^{M} S_i^{(m)}.
\]
Это позволяет объединять мнения нескольких агентов, эвристик или языковых моделей в одно итоговое решение.

\section{Байесовская обработка неопределённости}

Чтобы учесть неполные или шумные данные, система моделирует признаки кандидатов как случайные величины:
\[
q_i \sim \mathcal{D}_q, \qquad r_i \sim \mathcal{D}_r, \qquad s_i \sim \mathcal{D}_s.
\]
При наличии наблюдаемых данных $D$ байесовское обновление записывается как:
\[
p(\theta \mid D) \propto p(D \mid \theta)p(\theta).
\]

Этот слой особенно важен в задачах здоровья и покупок, где необходимо различать надёжные данные и единичные аномалии.

\section{Определение мультеверта}

Под \emph{мультевертом} понимается субъектно-центрированная вселенная решений:
\[
\mathcal{M}(u) = \{E_u, H_u, B_u, P_u, R_u\},
\]
где:
\begin{itemize}[leftmargin=1.2em]
    \item $E_u$ — экономическое подпространство,
    \item $H_u$ — медицинское подпространство,
    \item $B_u$ — поведенческое и бытовое подпространство,
    \item $P_u$ — прогностическое подпространство,
    \item $R_u$ — подпространство рисков и устойчивости.
\end{itemize}

Все вычисления в системе интерпретируются относительно $\mathcal{M}(u)$, а не относительно абстрактной глобальной цели. В этом и заключается ключевое отличие Angel-in-Pocket от обычных рекомендательных систем.

\section{Сценарии применения}

\subsection{Покупки}
Система сравнивает товары по цене, качеству, совместимости, доставке и доверию к продавцу. Она удаляет подозрительные предложения и ранжирует оставшиеся по долгосрочной ценности.

\subsection{Здоровье}
Система хранит и сравнивает анализы во времени, выявляет тренды и объясняет их простым языком. Она подчёркивает элементы, требующие медицинского внимания, но не заменяет врача.

\subsection{Образ жизни}
Система рекомендует корректировки сна, питания, физической активности и распределения времени. Она может прогнозировать, как изменения привычек повлияют на устойчивость в будущем.

\subsection{Прогнозирование}
На основе текущего состояния субъекта система симулирует траектории сценариев:
\[
x_{t+1} = F(x_t, a_t, \eta_t),
\]
где $a_t$ — рекомендованное действие, а $\eta_t$ — шум среды. Система оценивает несколько контрфактических сценариев:
\begin{itemize}[leftmargin=1.2em]
    \item что будет, если пользователь сохранит текущие привычки,
    \item что будет при улучшении диеты,
    \item что будет при росте стресса,
    \item что будет при ухудшении сна.
\end{itemize}

\section{Практическая реализация}

Практическая реализация может включать:
\begin{itemize}[leftmargin=1.2em]
    \item \texttt{backend/} для сервисов FastAPI,
    \item \texttt{frontend/} для веб- или десктопного интерфейса,
    \item \texttt{core/} для модулей профиля субъекта, скоринга и устойчивости,
    \item \texttt{health/}, \texttt{shopping/} и \texttt{services/} как адаптеры,
    \item \texttt{nullification/} и \texttt{swarm/} как отдельные модули,
    \item локальное хранилище памяти или зашифрованное пользовательское хранилище.
\end{itemize}

Систему можно расширить до настольного приложения, мобильного клиента или локально-ориентированного развертывания. Минимальный рабочий прототип требует лишь:
\begin{enumerate}[leftmargin=1.2em]
    \item ввода профиля субъекта,
    \item извлечения кандидатов,
    \item обнуления,
    \item ранжирования,
    \item генерации объяснения.
\end{enumerate}

\section{Обсуждение}

Предлагаемая архитектура поднимает важные вопросы конфиденциальности, доверия и автономии. Поскольку система знает интимные детали о субъекте, она должна проектироваться с сильной защитой данных и прозрачной логикой рекомендаций. Целью является не наблюдение, а расширение возможностей пользователя.

Ещё один важный аспект — эпистемическая скромность. Система не должна заявлять о точности медицинских или жизненных прогнозов. Вместо этого она должна сообщать вероятности, интервалы риска и уровни уверенности. Хороший Angel-in-Pocket не диктует жизнь; он помогает в ней ориентироваться.

\section{Заключение}

Angel-in-Pocket представляет собой субъектно-центрированную мультеверс-архитектуру для персонального ИИ-сопровождения. Её основная идея — перейти от поведения универсального ассистента к непрерывной, устойчивой и персонализированной системе принятия решений. Объединяя байесовскую обработку неопределённости, иерархическую устойчивость, обнуление низкокачественных сигналов и swarm-агрегацию, система может поддерживать практические сценарии в области покупок, здоровья и планирования жизни.

Данная архитектура подходит для поэтапной разработки и может служить фундаментом для реального личного ИИ-компаньона. В дальнейшем следует сосредоточиться на приватном хранении памяти, безопасном медицинском прогнозировании и объяснимости решений.

\section*{Благодарности}

Работа вдохновлена более широким исследовательским экосистемным контекстом GRA и идеей субъектно-центрированного ИИ-компаньона.

\begin{thebibliography}{9}

\bibitem{gra-core}
GRA-Core: Unified Hierarchical Stability Library. Концепция репозитория и архитектурного ядра.

\bibitem{bayesian}
Байесовские методы GRA для ранжирования и прогнозирования с учётом неопределённости.

\bibitem{swarm}
GRA-LLM-Swarm Constructs и GRA-Swarm Optimization.

\bibitem{nullification}
Методологии GRA Resonance Nullification и Chiral Nullification.

\bibitem{subject}
Концепции SubjectSwap и субъектно-центрированной оркестрации.

\end{thebibliography}

\end{document}